\section{Limitations and Trade-offs}
\tempo{} is designed as a lightweight library-level runtime. This choice comes
with several limitations that define the current scope:
\begin{itemize}
\item \textbf{No static guarantees.} Determinism and correct use of primitives
  (e.g., single emission on event signals) are enforced at runtime rather than
  by a type system or compiler checks.
\item \textbf{Runtime errors for misuse.} Multiple emissions on an event signal
  raise an exception; misuse must be caught by tests or careful discipline.
\item \textbf{Effects require OCaml 5.3.} The approach relies on algebraic
  effects and deep handlers, which are not available in earlier versions.
\item \textbf{Library-level scheduling only.} \tempo{} does not integrate with a
  language-level optimizer, so certain static analyses or compile-time
  optimizations (as in ReactiveML) are unavailable.
\item \textbf{Low-level API risks.} The \texttt{Tempo.Low\_level} module exposes
  primitives that can break high-level invariants if used incorrectly.
\item \textbf{Determinism vs evaluation order.} Reference-based communication
  can make behavior depend on evaluation order within an instant, which is not
  fixed by the semantics; mitigating this requires static analyses (see
  Section~\ref{sec:future-work}).
\item \textbf{Purity assumptions.} Aggregate signals assume order-insensitive
  combine functions; if combine is not associative/commutative, aggregate
  results may depend on intra-instant scheduling.
\item \textbf{Logical parallelism only.} Concurrency is logical and single-core;
  enabling physical parallelism would require additional determinism controls
  (see Section~\ref{sec:future-work}).
\item \textbf{Non-guaranteed behaviors.} Programs relying on shared mutable
  state across parallel tasks can observe different results even with identical
  inputs, because the intra-instant order is unspecified.
\end{itemize}
